Separator messages, cost shifting, and optimization over a tree are
established ideas. In finite-state graphical models, agreement on trees
leads to linear-programming and reparameterization interpretations
\citep{WainwrightEtAl2005}; AND/OR search exploits separator-conditioned
subproblems and cached states \citep{MarinescuDechter2009}.
Weighted constraint satisfaction combines tree decomposition with local
consistency and optimized cost shifts
\citep{deGivrySchiexVerfaillie2006,CooperDeGivrySchiex2007}.
For continuous pairwise models, Wald and Globerson
\citep{WaldGloberson2014} impose agreement on selected expectations and
derive a dual consisting of cost-shifted local minima. Their exactness
result for convex-decomposable pairwise models has a different hypothesis
from our growth-based construction. The restricted-split duality in
\cref{thm:consistency-duality} belongs to this line of work; we give its
compact minimax proof for the function classes used here. Unrestricted
tree gluing is also an established mechanism in sparse polynomial
optimization \citep{Lasserre2006}.

Positive-margin splitting under running intersection is central to sparse
polynomial representations. Grimm, Netzer, and Schweighofer
\citep{GrimmNetzerSchweighofer2007} construct such splits qualitatively;
Korda, Magron, and R\'ios-Zertuche
\citep{KordaMagronRiosZertuche2025} quantify polynomial approximation of
separator functions. Local polynomial shifts also appear in tightness
criteria for sparse moment relaxations \citep{NieQuTangZhang2026}.
Factor-two approximation bounds have precedents in approximate linear
programming \citep{deFariasVanRoy2003}, and moment matching gives a
zero-width approximation duality \citep{HanJiaoWeissman2018}.
Approximation of interval-valued bands through oriented distance is
studied by Ginchev and Hoffmann \citep{GinchevHoffmann2002}.
For a real interval band, their objective is
$\sup_s\max\{L(s)-\chi(s),\chi(s)-U(s)\}$. When the band pinches,
this is the uniform distance from $\chi$ to the band. Thus this
approximation objective is established.
We do not claim band approximation, qualitative splitting, or the
factor-two approximation principle as new. Our band results identify the
further equality between the optimized restricted dynamic-programming
gap and twice this distance, and distinguish separate-edge approximation
from one jointly exact tree split. The graded split and scalar covering
proof use the margin left in these bands.

Continuous decomposition methods provide closer algorithmic comparisons.
Berenguel et al. \citep{BerenguelEtAl2013} study interval branch-and-bound
for additively separable functions with common variables. Their
shared-variable pruning rule combines interval lower bounds from retained
subboxes whose common-variable intervals overlap the current interval
\citep[Proposition~1 and Theorem~2]{BerenguelEtAl2013}.
This supplies a direct antecedent for compatible-box decomposition bounds.
Our construction communicates affine separator minorants on a tree and
bounds the size and work of the resulting global certificate.
Bienstock and Mu\~noz \citep{BienstockMunoz2018} construct
bounded-treewidth linear-programming formulations for polynomial
optimization, with explicit dependence on inverse accuracy.
Zhang and Sun \citep{ZhangSun2022} give accuracy-dependent iteration
bounds for multistage nonconvex decomposition under their state-covering
and oracle assumptions. MUSE-BB \citep{LangiuEtAl2025} uses decomposable
scenario subproblems and multisection branching in one spatial search
tree. These methods establish the value of decomposition, but their
stated guarantees differ from our final affine-certificate size and
construction count under point growth.
Robertson, Cheng, and Scott \citep{RobertsonChengScott2025} analyze
convergence orders of reduced-space value-function relaxations for
nonconvex stochastic programs.

Adaptive grids are likewise established. Munos and Moore
\citep{MunosMoore2002} refine state-space cells and rebuild discretized
dynamic programs in optimal control. Nagarajan et al.
\citep{NagarajanEtAl2019} combine adaptive multivariate partitioning,
piecewise relaxations, and bound tightening for global optimization.
Our contribution is not adaptivity itself. It is the aggregate disagreement
estimate that makes a fresh graded partition contract around a point
obtained from an affine separator certificate, on an arbitrary supplied
tree. The logarithmic accuracy dependence relies on global quadratic
growth and the displayed structural and conditioning parameters.
Logarithmic accuracy dependence also occurs in optimistic optimization
under suitable local-growth and partition assumptions
\citep{Munos2011OptimisticOptimization}. That guarantee concerns simple
regret. Certified Lipschitz optimization has instance-dependent value-oracle
bounds \citep{BachocCesariGerchinovitz2021}. Our certificate and work counts
use different information: bag gradients, local lower models, and local
convex minimizations. These oracle models should not be compared as if
they supplied the same operations.

Inexact dynamic programming and inexact first-order methods have their own
error-propagation theory \citep{Munos2005,DevolderGlineurNesterov2014}.
Guigues develops deterministic and stochastic inexact decomposition from
approximate primal and dual subproblem solves
\citep{Guigues2017InexactCuts,Guigues2020InexactSDDP}. These related analyses
track stage errors and permit increasingly accurate solves; they already
establish inexact decomposition as a method.
Reliable rational lower bounds from numerical linear programming are
developed by Steffy and Wolter \citep{SteffyWolter2013}, and
limited-precision LP oracles can be used to recover exact solutions
\citep{GleixnerSteffy2020}. We use a specific certified contract: an exact
local feasible point, a valid lower bound, an error budget proportional to
squared cell width, and an aggregate separator-slope error. Its residual
identity charges one local error per selected bag. For polynomial box
objectives, a separate affine Taylor construction eliminates the general
convex oracle: every local affine problem is solved at interval endpoints,
and the resulting rational certificate has a compact representation.
For fixed degree, bag dimension, occurrence, and numerical conditioning,
the bit cost is polynomial in input length and precision. It is not
uniformly polynomial in the curvature-to-growth ratio. This realization
selects a particular aggregate bag model; it does not implement every
possible factorwise relaxation.

Backward adjoints are standard in nonlinear optimal control
\citep{Hager2000}. The dynamics extension uses them to cancel the
first-order objective error of forward feasibility repair in a separator
certificate. The cancellation holds at ambient-box centers as well as
feasible centers. Contracting graph maps then permit short certified
enclosures and repeated center resets. The assumptions include full-box
invariance, feasible point growth, and sound derivative bounds; this is a
specific constrained extension rather than a theorem for arbitrary
nonlinear constraints.

To the best of our knowledge, the results reviewed above do not give the
combination proved here: construction of affine separator certificates
without a known optimizer on arbitrary tree decompositions, the final
logarithmic partition count and explicit incidence-based oracle count
under point growth, certified local error budgets for that construction,
and the stated rational implementations. The claims concern this
combination and these precise contracts. They do not assert novelty of
tree dynamic programming, cost shifting, adaptive partitioning, adjoints,
or exact rational linear programming.
